
\documentclass{article}
\usepackage{colortbl}
\usepackage{makecell}
\usepackage{multirow}
\usepackage{supertabular}

\begin{document}

\newcounter{utterance}

\centering \large Interaction Transcript for game `eqbench', experiment `eq\_bench', episode 55 with reflexion{-}Qwen3.5{-}27B.
\vspace{24pt}

{ \footnotesize  \setcounter{utterance}{1}
\setlength{\tabcolsep}{0pt}
\begin{supertabular}{c@{$\;$}|p{.15\linewidth}@{}p{.15\linewidth}p{.15\linewidth}p{.15\linewidth}p{.15\linewidth}p{.15\linewidth}}
    \# & \multicolumn{2}{c}{Player} && \multicolumn{2}{c}{Game Master} \\
    \hline

    \theutterance \stepcounter{utterance}  
    & & & \multicolumn{4}{p{0.6\linewidth}}{
        \cellcolor[rgb]{0.9,0.9,0.9}{
            \makecell[{{p{\linewidth}}}]{
                \texttt{\tiny{[P1$\langle$GM]}}
                \\ 
\texttt{Your task is to predict the likely emotional responses of a character in this dialogue:} \\
\\ 
\texttt{Reggie: Marcus, I've been thinking about Mom's birthday party. I've got some ideas.} \\
\texttt{Marcus: Oh, I bet you do. Another corporate gala, Reggie? } \\
\texttt{Reggie: Funny. I was actually thinking of something more... intimate. } \\
\texttt{Marcus: Intimate? You? Now, that's a surprise.} \\
\texttt{Reggie: Yes, well, I can be full of surprises. } \\
\texttt{Marcus: I'll believe it when I see it. } \\
\texttt{Reggie: What's that supposed to mean? } \\
\texttt{Marcus: Nothing. Just that your idea of surprise usually involves a powerpoint presentation and a five{-}year plan.} \\
\texttt{Reggie: And your idea of a surprise is forgetting to show up. } \\
\texttt{Marcus: Low blow, sis. } \\
\texttt{{[}End dialogue{]}} \\
\\ 
\texttt{At the end of this dialogue, Marcus would feel...} \\
\texttt{Defeated} \\
\texttt{Provoked} \\
\texttt{Apologetic} \\
\texttt{Indifferent} \\
\\ 
\texttt{Give each of these possible emotions a score from 0{-}10 for the relative intensity that they are likely to be feeling each.} \\
\\ 
\texttt{You must output in the following format, including headings (of course, you should give your own scores), with no additional commentary:} \\
\\ 
\texttt{Defeated: <score>} \\
\texttt{Provoked: <score>} \\
\texttt{Apologetic: <score>} \\
\texttt{Indifferent: <score>} \\
\\ 
\\ 
\texttt{{[}End of answer{]}} \\
\\ 
\texttt{Remember: zero is a valid score, meaning they are likely not feeling that emotion. You must score at least one emotion > 0.} \\
\\ 
\texttt{Your answer:} \\
            }
        }
    }
     \\ \\

    \theutterance \stepcounter{utterance}  
    & \multicolumn{4}{p{0.6\linewidth}}{
        \cellcolor[rgb]{0.561,0.737,0.561}{
            \makecell[{{p{\linewidth}}}]{
                \texttt{\tiny{[P1$\rangle$GM]}}
                \texttt{Defeated: 2} \\
\texttt{Provoked: 8} \\
\texttt{Apologetic: 1} \\
\texttt{Indifferent: 0} \\
            }
        }
    }
    & & \\ \\

    \theutterance \stepcounter{utterance}  
    & & & \multicolumn{2}{p{0.3\linewidth}}{
        \cellcolor[rgb]{0.9,0.9,0.9}{
            \makecell[{{p{\linewidth}}}]{
                \texttt{\tiny{[GM$|$GM]}}
                \texttt{Defeated: 2} \\
\texttt{Provoked: 8} \\
\texttt{Apologetic: 1} \\
\texttt{Indifferent: 0} \\
            }
        }
    }
    & & \\ \\

    \theutterance \stepcounter{utterance}  
    & & & \multicolumn{2}{p{0.3\linewidth}}{
        \cellcolor[rgb]{0.9,0.9,0.9}{
            \makecell[{{p{\linewidth}}}]{
                \texttt{\tiny{[GM$|$GM]}}
                \texttt{Defeated: 0} \\
\texttt{Provoked: 7} \\
\texttt{Apologetic: 0} \\
\texttt{Indifferent: 0} \\
            }
        }
    }
    & & \\ \\

    \theutterance \stepcounter{utterance}  
    & & & \multicolumn{2}{p{0.3\linewidth}}{
        \cellcolor[rgb]{0.9,0.9,0.9}{
            \makecell[{{p{\linewidth}}}]{
                \texttt{\tiny{[GM$|$GM]}}
                \texttt{game\_result = LOSE} \\
            }
        }
    }
    & & \\ \\

    \theutterance \stepcounter{utterance}  
    & & \multicolumn{4}{p{0.6\linewidth}}{
        \cellcolor[rgb]{0.95,0.95,0.95}{
            \makecell[{{p{\linewidth}}}]{
                \texttt{\tiny{[P1$\leftrightarrow$P]}}
                \texttt{Post{-}episode Reflexion (not sent to GM):} \\
\texttt{New plan: My previous attempt failed because I assigned a high score to "Provoked" (8) while ignoring the nuance of the final exchange. Marcus's line "Low blow, sis" is a defensive reaction to a personal attack, but the context of the entire dialogue shows a pattern of mutual bickering where Marcus consistently deflects with sarcasm ("I'll believe it when I see it") rather than genuine anger. By scoring "Provoked" so high, I likely missed the target state which probably required a more balanced assessment or a different primary emotion like "Defeated" (due to the "Low blow" admission) or "Indifferent" (given his dismissive tone). Next time, I will analyze the final line more critically: "Low blow" acknowledges a hit but doesn't necessarily mean he is furious; it often signals a grudging admission of defeat in an argument. I will lower the score for "Provoked" and increase the score for "Defeated" to reflect that his sarcasm has been punctured, or consider "Indifferent" if the target implies he doesn't care about the fight. I will ensure the scores reflect the specific dynamic of a sibling argument where the "winner" is often the one who lands the final personal jab, leaving the other slightly defeated rather than enraged.} \\
            }
        }
    }
    & \\ \\

\end{supertabular}
}

\end{document}
